\subsection{M-flat modules}

DEFINITION 1. Let A be a ring, E a right A-module and M a left A-module. E is said to be flat for M (or M-flat) if, for every left A-module M' and every injective homomorphism \(v: M' \to M\), the homomorphism \(1_{E} \otimes v\) : E \(\otimes_{A} M' \to E \otimes_{A} M\) is injective.

For any right A-module N, the notion of an N-flat left module is defined similarly. To say that a right A-module E is flat for a left A-module M is equivalent to saying that E, considered as a left \(A^{0}\) -module (recall that \(A^{0}\) denotes the opposite ring to A), is flat for the right \(A^{0}\) -module M.

LEMMA 3. For a right A-module E to be M-flat, it is necessary and sufficient that for every finitely generated submodule M' of M the canonical homomorphism

\[
\mathbf {1} _ {\mathbf {E}} \otimes j: \mathbf {E} \otimes_ {\mathbf {A}} \mathbf {M} ^ {\prime} \rightarrow \mathbf {E} \otimes_ {\mathbf {A}} \mathbf {M}
\]

(j being the canonical injection \(\mathbf{M}'\to \mathbf{M}\) ) be injective.

Suppose that this condition holds and let N be any submodule of M. Suppose that the canonical image in \(E \otimes_{A} M\) of an element

\[
z = \sum_ {i} x _ {i} \otimes y _ {i} \in \mathbf {E} \otimes_ {\mathbf {A}} \mathbf {N} (x _ {i} \in \mathbf {E}, y _ {i} \in \mathbf {N})
\]

is zero and let \(M'\) be the finitely generated submodule of N generated by the \(y_{i}\) ; as by hypothesis the composite mapping \(E \otimes_{A} M' \to E \otimes_{A} N \to E \otimes_{A} M\) is injective, the sum \(z' = \sum_{i} x_{i} \otimes y_{i}\) , considered as an element of \(E \otimes_{A} M'\) , is zero. As z is the image of \(z'\) , we have also z = 0, whence the lemma.

LEMMA 4. Let E be a right A-module and M a left A-module such that E is M-flat. If N is either a submodule or a quotient module of M, then E is N-flat.

The case in which N is a submodule is easy, as, if \(N'\) is a submodule of N, the composite homomorphism

\[
\mathrm{E} \otimes_ {\mathrm{A}} \mathrm{N} ^ {\prime} \rightarrow \mathrm{E} \otimes_ {\mathrm{A}} \mathrm{N} \rightarrow \mathrm{E} \otimes_ {\mathrm{A}} \mathrm{M}
\]

is injective, hence so is \(\mathbf{E} \otimes_{\mathbf{A}} \mathbf{N}' \to \mathbf{E} \otimes_{\mathbf{A}} \mathbf{N}\) . Suppose then that \(\mathbf{N}\) is a quotient module of \(\mathbf{M}\) , that is there exists an exact sequence \(0 \to \mathbf{R} \xrightarrow{i} \mathbf{M} \xrightarrow{p} \mathbf{N} \to 0\) . Let \(\mathbf{N}'\) be a submodule of \(\mathbf{N}\) and \(\mathbf{M}' = \bar{p}^{-1}(\mathbf{N}')\) . Let \(i'\) denote the mapping of \(\mathbf{R}\) to \(\mathbf{M}'\) with the same graph as \(i, p'\) the surjection \(\mathbf{M}' \to \mathbf{N}'\) with the same graph as the restriction of \(p\) to \(\mathbf{M}'\) , \(r\) the identity mapping of \(\mathbf{R}\) to \(\mathbf{R}\) , \(m\) the canonical injection \(\mathbf{M}' \to \mathbf{M}\) and \(n\) the canonical injection \(\mathbf{N}' \to \mathbf{N}\) . The diagram

\[
\begin{array}{c} 0 \longrightarrow \mathrm{R} \xrightarrow {i ^ {\prime}} \mathrm{M} ^ {\prime} \xrightarrow {p ^ {\prime}} \mathrm{N} ^ {\prime} \longrightarrow 0 \\ r \Biggl \downarrow \qquad m \Biggl \downarrow \qquad n \Biggl \downarrow \\ 0 \longrightarrow \mathrm{R} \xrightarrow [ i ]{} \mathrm{M} \xrightarrow [ p ]{} \mathrm{N} \longrightarrow 0 \end{array}
\]

is commutative and its rows are exact.

To simplify the writing, we set \(\mathrm{T}(\mathbf{Q}) = \mathrm{E} \otimes_{\mathbf{A}} \mathbf{Q}\) for every left A-module \(\mathbf{Q}\) and \(\mathrm{T}(v) = 1_{\mathbf{E}} \otimes v\) for every left A-module homomorphism \(v\) . The diagram

\[
\begin{array}{l} \mathrm{T} (\mathrm{R}) \xrightarrow {\mathrm{T} (i ^ {\prime})} \mathrm{T} (\mathrm{M} ^ {\prime}) \xrightarrow {\mathrm{T} (p ^ {\prime})} \mathrm{T} (\mathrm{N} ^ {\prime}) \longrightarrow 0 \\ \mathrm{T} (r) \Bigg \downarrow \qquad \qquad \qquad \qquad \qquad \qquad \qquad \qquad \qquad \qquad \qquad \qquad \qquad \qquad \qquad \qquad \qquad \qquad \qquad \qquad \qquad \qquad \qquad \qquad \qquad \qquad \qquad \qquad \qquad \qquad \qquad \qquad \qquad \qquad \\ \mathrm{T} (\mathrm{R}) \xrightarrow [ \mathrm{T} (i) ]{} \mathrm{T} (\mathrm{M}) \xrightarrow [ \mathrm{T} (p) ]{} \mathrm{T} (\mathrm{N}) \longrightarrow 0 \end{array}
\]

is commutative and its rows are exact by Lemma 1 of no. 1. Moreover, since E is M-flat, the homomorphism T(m) is injective. As T(r) and T(p') are surjective, it follows from § 1, no. 4, Corollary 2 to Proposition 2 that T(n) is injective, which proves the lemma.

LEMMA 5. Let \((\mathbf{M}_{\iota})_{\iota \in \mathbf{I}}\) be a family of left A-modules, \(\mathbf{M} = \bigoplus_{\iota \in \mathbf{I}} \mathbf{M}_{\iota}\) their direct sum and E a right A-module. If, for all \(\iota \in \mathbf{I}\) , E is flat for \(\mathbf{M}_{\iota}\) , then E is flat for M.

\begin{enumerate}
\def\labelenumi{(\alph{enumi})}
\tightlist
\item
  Suppose first that \(I = \{1, 2\}\) , and let \(M'\) be a submodule of \(M = M_{1} \oplus M_{2}\) , \(M_{1}\) and \(M_{2}\) being canonically identified with submodules of M. Denote by \(M_{1}'\) the intersection \(M' \cap M_{1}\) and by \(M_{2}'\) the image of \(M'\) in \(M_{2}\) under the canonical projection p of M onto \(M_{2}\) . We have a diagram
\end{enumerate}

\[
\begin{array}{c} 0 \longrightarrow \mathrm{M} _ {1} ^ {\prime} \xrightarrow {i ^ {\prime}} \mathrm{M} ^ {\prime} \xrightarrow {p ^ {\prime}} \mathrm{M} _ {2} ^ {\prime} \longrightarrow 0 \\ v _ {1} \Biggl \downarrow \qquad \qquad \qquad \qquad \qquad \qquad \qquad \qquad \qquad \qquad \qquad \qquad \qquad v _ {2} \Biggl \downarrow \\ 0 \longrightarrow \mathrm{M} _ {1} \xrightarrow [ i ]{} \mathrm{M} \xrightarrow [ p ]{} \mathrm{M} _ {2} \longrightarrow 0 \end{array}
\]

where \(v_{1}, v, v_{2}, i, i'\) are the canonical injections and \(p'\) is the mapping with the same graph as the restriction of p to M', which is surjective. It is immediately verified that this diagram is commutative and that its rows are exact. With T(Q) and T(v) used in the same sense as in the proof of Lemma 4, we have a commutative diagram

\[
\begin{array}{c} \mathrm{T} (\mathrm{M} _ {1} ^ {\prime}) \xrightarrow {\mathrm{T} (i ^ {\prime})} \mathrm{T} (\mathrm{M} ^ {\prime}) \xrightarrow {\mathrm{T} (p ^ {\prime})} \mathrm{T} (\mathrm{M} _ {2} ^ {\prime}) \\ \mathrm{T} (v _ {1}) \Biggl \downarrow \qquad \qquad \qquad \qquad \qquad \qquad \qquad \qquad \qquad \qquad \qquad \qquad \qquad \qquad \qquad \qquad \qquad \qquad \qquad \qquad \qquad \qquad \qquad \qquad \qquad \qquad \qquad \qquad \qquad \qquad \qquad \qquad \qquad \qquad \\ \mathrm{T} (\mathrm{M} _ {1}) \xrightarrow [ \mathrm{T} (i) ]{} \mathrm{T} (\mathrm{M}) \xrightarrow [ \mathrm{T} (p) ]{} \mathrm{T} (\mathrm{M} _ {2}) \end{array}
\]

By Lemma 1 of no. 1 the two rows of this diagram are exact; as E is flat for \(M_{1}\) and \(M_{2}\) , \(T(v_{1})\) and \(T(v_{2})\) are injective; moreover, by Lemma 2 of no. 1, \(T(i)\) is injective. Corollary 2 to Proposition 2 of § 1, no. 4 then shows that \(T(v)\) is injective and consequently E is M-flat.

\begin{enumerate}
\def\labelenumi{(\alph{enumi})}
\setcounter{enumi}{1}
\item
  If I is a finite set with \(n\) elements, the lemma follows by induction on \(n\) using (a).
\item
  In the general case let \(\mathbf{M}'\) be a finitely generated submodule of \(\mathbf{M}\) . Then there exists a finite subset \(J\) of the indexing set \(I\) such that \(\mathbf{M}'\) is contained in the direct sum \(\mathbf{M}_J = \bigoplus_{i \in J} \mathbf{M}_i\) . By (b) \(E\) is flat for \(\mathbf{M}_J\) ; the canonical homomorphism \(\mathbf{E} \otimes_{\mathbf{A}} \mathbf{M}' \to \mathbf{E} \otimes_{\mathbf{A}} \mathbf{M}_J\) is therefore injective. On the other hand, as \(\mathbf{M}_J\) is a direct factor of \(\mathbf{M}\) , the canonical homomorphism \(\mathbf{E} \otimes_{\mathbf{A}} \mathbf{M}_J \to \mathbf{E} \otimes_{\mathbf{A}} \mathbf{M}\) is injective (no. 1, Lemma 2). Taking the composition, it follows that \(\mathbf{E} \otimes_{\mathbf{A}} \mathbf{M}_J \to \mathbf{E} \otimes_{\mathbf{A}} \mathbf{M}\) is injective and \(E\) is flat for \(M\) by Lemma 3.
\end{enumerate}

